
% ╔═══════════════════════╗
% ║  Theorems and proofs  ║
% ╚═══════════════════════╝

% Theorem
% -------------------------------------------------------------------------------------
\newtcolorbox{theorem}[2][] {
    %
    % Layout
    enhanced, breakable, before skip=10pt, after skip=10pt,
    before upper = {\vspace{3mm}},
    %
    % Box
    boxrule = 2pt,
    colframe = MAIN, colback = BG,
    sharp corners,
    drop shadow,
    %
    % Title
    attach boxed title to top left = {
        xshift=1.3cm, yshift=-4mm
    },
    boxed title style = {
        boxrule = 2pt,
        colframe = MAIN, colback = BG,
        sharp corners,
        underlay = {
            \begin{scope}
                \clip ([xshift=2pt, yshift=9.4pt]interior.south east) rectangle ([xshift=4pt]interior.north east);
                \filldraw [white, opacity=0, drop shadow] (interior.south west) rectangle ([xshift=2pt, yshift=2pt]interior.north east);
            \end{scope} 
        },
    },
    boxed title size = title,
    coltitle = MAIN,
    fonttitle = \bfseries,
    title={#2},
    %
    % Text
    colupper = TEXT,
    %
    % Border fixes for multiple pages
    overlay first={
        \draw[color=BG, line width = 4pt] ([yshift=-1pt]frame.south east)--([yshift=-1pt]frame.south west);
    },
    overlay middle={
        \draw[color=BG, line width = 2pt] (frame.south east)--(frame.south west);
        \draw[color=BG, line width=2pt] (frame.north east)--(frame.north west);
    },
    overlay last={
        \draw[color=BG, line width=2pt] (frame.north east)--(frame.north west);
    },
    %
    % User input parameters
    #1
}
%
% Usage:
% \begin{theorem} [tcolorbox params] {Theorem name}
%     Formulation
% \end{theorem}
% -------------------------------------------------------------------------------------


% Theorem proof
% -------------------------------------------------------------------------------------
\newtcolorbox{tproof}[1][] {
    %
    % Layout
    enhanced, breakable, before skip=10pt, after skip=10pt,
    after upper = {\vspace{-2mm}},
    %
    % Box
    opacityback = 0,
    frame hidden,
    %
    % CATS !!!!!!!!!!!!!
    % watermark graphics = cats.png,
    % watermark opacity = .1,
    % watermark overzoom = 1,
    %
    % Proof title
    title = {Доказательство.},
    attach title to upper = {\ },
    coltitle = MAIN,
    fonttitle = \bfseries,
    %
    % Text
    colupper = TEXT,
    %
    % Line
    borderline west = {2pt}{0pt}{MAIN},
    %
    % Square
    underlay unbroken and last = {
        \fill[color=MAIN] ([yshift=1.5mm]frame.south east) rectangle ++(-3mm, 3mm);
    },
    %
    % User input parameters
    #1
}
%
% Usage:
% \begin{tproof} [tcolorbox params] 
%    Proof 
% \end{tproof}
% -------------------------------------------------------------------------------------


% Определение фиолетового цвета для рамок
\definecolor{EXAMP_COLOR}{RGB}{102, 51, 153} 

% Example Box
% -------------------------------------------------------------------------------------
\newtcolorbox{example}[2][] {
    enhanced, breakable, before skip=10pt, after skip=10pt,
    before upper = {\vspace{3mm}},
    boxrule = 2pt,
    colframe = EXAMP_COLOR, 
    colback = BG, % ИСПОЛЬЗУЕМ ВАШ ЦВЕТ ФОНА ВМЕСТО БЕЛОГО
    sharp corners,
    drop shadow,
    %
    attach boxed title to top left = {
        xshift=1.3cm, yshift=-4mm
    },
    boxed title style = {
        boxrule = 2pt,
        colframe = EXAMP_COLOR, 
        colback = BG, % И ЗДЕСЬ ТОЖЕ
        sharp corners,
        underlay = {
            \begin{scope}
                \clip ([xshift=2pt, yshift=9.4pt]interior.south east) rectangle ([xshift=4pt]interior.north east);
                \filldraw [white, opacity=0, drop shadow] (interior.south west) rectangle ([xshift=2pt, yshift=2pt]interior.north east);
            \end{scope} 
        },
    },
    boxed title size = title,
    coltitle = EXAMP_COLOR,
    fonttitle = \bfseries,
    title={Пример: #2},
    colupper = TEXT,
    %
    % Исправление границ для многостраничности (теперь прозрачно сливается с фоном)
    overlay first={
        \draw[color=BG, line width = 4pt] ([yshift=-1pt]frame.south east)--([yshift=-1pt]frame.south west);
    },
    overlay middle={
        \draw[color=BG, line width = 2pt] (frame.south east)--(frame.south west);
        \draw[color=BG, line width=2pt] (frame.north east)--(frame.north west);
    },
    overlay last={
        \draw[color=BG, line width=2pt] (frame.north east)--(frame.north west);
    },
    #1
}

% Решение (eproof)
\newtcolorbox{eproof}[1][] {
    enhanced, breakable, before skip=10pt, after skip=10pt,
    after upper = {\vspace{-2mm}},
    opacityback = 0, % Делаем фон прозрачным, чтобы просвечивал BG
    frame hidden,
    title = {Решение.},
    attach title to upper = {\ },
    coltitle = EXAMP_COLOR,
    fonttitle = \bfseries,
    colupper = TEXT,
    borderline west = {2pt}{0pt}{EXAMP_COLOR},
    underlay unbroken and last = {
        \fill[color=EXAMP_COLOR] ([yshift=1.5mm]frame.south east) rectangle ++(-3mm, 3mm);
    },
    #1
}